\chapter*{Реферат}

Отчёт состоит из 1 книги без приложений, включающей реферат, термины и определения, обозначения и сокращения, введение, 6 разделов, заключение, с количеством стр. -- \thelastpage, иллюстраций -- \total{figure}, табл. -- \total{table}, использованных источников -- \total{citesnum}.

Перечень ключевых слов:

дифференцируемый рендеринг, обратный рендеринг, функция расстояния со знаком, разреженное множество бриков, редистансинг, реконструкция освещения, гониометрическая диаграмма направленности, кэш световых путей, трассировка путей.

Основание для проведения работы:

Основанием для проведения работы является соглашение о предоставлении гранта в форме субсидии от 20 июня 2025 года №~139-15-2025-012, заключенного с Министерством экономического развития Российской Федерации в рамках реализации федерального проекта «Искусственный интеллект», входящего в состав национального проекта «Экономика данных и цифровая трансформация государства».

Объект исследования:

Методы и алгоритмы дифференцируемого рендеринга для реконструкции поверхностей и конфигурации освещения.

Цель работы:

Разработка новых методов и алгоритмов прямого и обратного рендеринга, достигающих многократного выигрыша по вычислительным затратам и затратам по памяти у существующих методов, а также обеспечение их доступности на широком спектре вычислительных устройств за счёт создания кроссплатформенного программного обеспечения на базе собственных методов и алгоритмов.

Методология оценки результатов:

Точность реконструкции оценивается по метрике PSNR между изображением, рассчитанным с восстановленными параметрами сцены, и эталонным (входным) изображением. Эффективность методов оценивается по времени работы на эталонном оборудовании (центральный процессор AMD Ryzen~9 7950X, графический процессор NVIDIA GeForce RTX~2080) и по потреблению памяти в сравнении с аналогами, выбранными по результатам обзоров литературы: для реконструкции поверхностей -- по времени рендеринга, итерации и всей реконструкции, пиковой памяти и размеру модели; для реконструкции освещения -- по времени вычисления градиента и времени достижения одинакового уровня PSNR.

Результаты и их новизна:

Разработан метод дифференцируемого рендеринга поверхностей, заданных разреженным множеством бриков функции расстояния со знаком, с адаптированными к нему алгоритмами поиска пересечений и точек релаксированной границы и двумя алгоритмами редистансинга. По сравнению с базовым методом Relaxed Boundary итерация реконструкции выполняется в 6,8--10,3 раза быстрее, представление на разрешении $512^3$ компактнее плотной сетки в 4,8--15 раз, а точность после 7 эпох выше в среднем на 7,2~дБ. Разработан метод дифференцируемого рендеринга для реконструкции источников освещения с гониометрической диаграммой направленности и пространственно неоднородных площадных источников по HDR-изображению с учётом глобального освещения, основанный на однократном построении кэша световых путей; он достигает того же PSNR, что и система Mitsuba~3, в 2,6--4,3 раза быстрее. Разработана система LightCadRenderCore, выполняющая реконструкцию источника по фотографии сцены, расчёт изображения сцены с восстановленным источником и экспорт источника в стандартный формат IES.

Внедрение:

Внедрение разработанных методов реконструкции источников освещения предполагается в рамках прикладного проекта с ООО ИНТИЛЕД \todo{уточнить}.

Эффект: амортизация роста требований к вычислительным ресурсам при применении методов и моделей ИИ в будущем ПО для решения различных задач работы с 3D-данными, их визуализации и реконструкции.

Область применения: компьютерная графика, инженерное проектирование, проектирование систем освещения, научная и инженерная визуализация.

Структура отчёта в соответствии с техническим заданием (раздел 3.1):

\begin{itemize}
    \item 3.1.7 Обзор литературы на тему эффективных методов дифференцируемого рендеринга для реконструкции или оптимизации конфигурации освещения: раздел~\ref{ch:light-review}.
    \item 3.1.8 Метод дифференцируемого рендеринга для реконструкции или оптимизации конфигурации освещения: раздел~\ref{ch:light-method}.
    \item 3.1.9 Результаты сравнения разработанного метода дифференцируемого рендеринга для реконструкции или оптимизации конфигурации освещения с аналогами: раздел~\ref{chap:light-comparison}.
    \item 3.1.10 Обзор литературы на тему эффективных методов дифференцируемого рендеринга поверхностей: раздел~\ref{ch:chap1}.
    \item 3.1.11 Метод дифференцируемого рендеринга поверхностей: раздел~\ref{ch:chap2}.
    \item 3.1.12 Результаты сравнения разработанного метода дифференцируемого рендеринга поверхностей с аналогами: раздел~\ref{ch:chap4}.
    \item 3.1.13 Система LightCadRenderCore, демонстрирующая возможности разработанных методов реконструкции освещения: подраздел~\ref{sec:lightcad} (описание) и подраздел~\ref{sec:lightcad-tests} (испытания).
\end{itemize}

\endinput
